\chapter{Implementation Details}
\label{app:implementation-details}

This appendix summarizes implementation details that are useful for
reproducing the experiments but would interrupt the main methodological
discussion. The repository was organized so that each task has its own
configuration files, dataset builder, model runners, saved model artifacts, and
result folders, while sharing common preprocessing, switch construction, and
evaluation utilities.

\section{Repository Structure}
\label{app:repository-structure}

The codebase separates common infrastructure from task-specific logic. Common
utilities live under \path{src/data/}, \path{src/switching/}, \path{src/utils/},
and \path{src/tuning/}. Task-specific implementations live under
\path{src/tasks/<task_name>/}. Experiment entry points are grouped in
\path{scripts/experiments/<task_name>/}, while analysis and comparison scripts
are grouped in \path{scripts/analysis/}.

\Cref{tab:appendix-repo-structure} gives the practical role of the main
repository folders. The table is intended as a reproducibility map: it shows
where configuration, immutable raw data, intermediate preprocessing outputs,
final model-ready datasets, fitted models, and result artifacts are stored.

\begin{table}[htbp]
  \centering
  \caption{Main repository folders used by the experimental pipeline.}
  \label{tab:appendix-repo-structure}
  \small
  \setlength{\tabcolsep}{4pt}
  \renewcommand{\arraystretch}{1.15}
  \begin{tabular}{
    @{}
    >{\raggedright\arraybackslash}p{0.30\textwidth}
    >{\raggedright\arraybackslash}p{0.62\textwidth}
    @{}
  }
    \toprule
    Folder & Purpose \\
    \midrule
    \path{configs/}
      & YAML configuration files for data preparation, model runs, grid searches, Perfect Switch construction, and switch comparisons. \\
    \path{data/raw/}
      & Original input files. These files are treated as read-only. \\
    \path{data/interim/}
      & Cleaned signals, candidate events, event windows, and intermediate reusable preprocessing artifacts. \\
    \path{data/processed/}
      & Final task-specific arrays and metadata used by model runners. \\
    \path{src/tasks/}
      & Task-scoped dataset, model, evaluation, plotting, and utility code. \\
    \path{src/switching/}
      & Model-agnostic Perfect Switch, switch conversion, post-processing, metrics, and plotting utilities. \\
    \path{scripts/experiments/}
      & Executable training, zero-shot evaluation, dataset-building, and grid-search scripts. \\
    \path{scripts/analysis/}
      & Data-analysis and post-training comparison scripts. \\
    \path{models/}
      & Saved checkpoints, fitted estimators, model configurations, and run metadata. \\
    \path{results/}
      & Predictions, metrics, tables, figures, comparison summaries, and result indexes. \\
    \bottomrule
  \end{tabular}
\end{table}

The main task namespaces are \path{autoregressive},
\path{current_level_persistence}, \path{long_fade_detection}, and
\path{survival_persistence}. This separation avoids storing task-specific
models or outputs under a generic folder and makes it possible to trace every
result back to its task formulation.

\section{Dataset Artifact Schema}
\label{app:dataset-artifact-schema}

Every processed dataset is stored together with metadata describing the
selection identifier, split definition, context length, threshold assumptions,
and preprocessing fingerprint. The canonical split is development versus
external holdout: training and validation are drawn from development sources,
whereas \path{fc-uplink-fade.csv} is used only as the final test source.

\Cref{tab:appendix-dataset-artifacts} summarizes the arrays and metadata saved
by each task-specific dataset builder. The table separates numerical arrays
from metadata because the metadata is required to align predictions back to
timestamps, events, source datasets, and switch-evaluation windows.

\begin{table}[htbp]
  \centering
  \caption{Main processed dataset artifacts by task.}
  \label{tab:appendix-dataset-artifacts}
  \small
  \setlength{\tabcolsep}{3pt}
  \renewcommand{\arraystretch}{1.15}
  \begin{tabular}{
    @{}
    >{\raggedright\arraybackslash}p{0.23\textwidth}
    >{\raggedright\arraybackslash}p{0.33\textwidth}
    >{\raggedright\arraybackslash}p{0.34\textwidth}
    @{}
  }
    \toprule
    Task & Array artifacts & Metadata artifacts \\
    \midrule
    Autoregressive forecasting
      & \path{train.npz}, \path{val.npz}, and \path{test.npz} with raw and context-standardized past windows and future trajectories.
      & Split metadata parquet files with window identifiers, event identifiers, source dataset names, timestamps, quality flags, and scaling parameters. \\
    Current-level persistence
      & Split NPZ files with raw context, relative-to-current context, scalar context features, and persistence duration targets.
      & Metadata with target duration, censoring flags, acquisition segment identifiers, scalar feature names, and train-only scaler parameters. \\
    Long-fade detection
      & Split NPZ files with sequence context, relative-to-threshold context, scalar context, and binary long-fade labels.
      & Metadata with grouped-event identifiers, event labels, timestamps, dataset identifiers, and quality flags. \\
    Survival persistence
      & Split NPZ files with past contexts, scalar features, residual-duration labels, and AFT lower/upper bounds.
      & Metadata with event identifiers, censoring indicators, acquisition segment identifiers, observed durations, and survival target bounds. \\
    \bottomrule
  \end{tabular}
\end{table}

The split metadata is as important as the numerical arrays. It preserves the
alignment between model inputs, event timestamps, source files, event
identifiers, and later switch comparisons. Model scripts do not reconstruct
splits independently; they load these artifacts and validate their fingerprints
where necessary.

\section{Run and Result Artifact Schema}
\label{app:run-result-artifact-schema}

Each trained or evaluated model writes a run folder under
\path{results/runs/<task_name>/} and, when applicable, a model folder under
\path{models/<task_name>/}. Run identifiers encode the task, model family,
variant, selection identifier, and relevant setup. A typical run folder
contains:

\begin{itemize}
  \item \path{predictions/}: validation and/or test predictions in parquet
  format;
  \item \path{metrics/}: native task metrics such as forecast errors,
  classification scores, duration errors, or survival metrics;
  \item \path{figures/}: diagnostic plots generated by the run;
  \item \path{tables/}: compact run summaries;
  \item \path{metadata.yaml}: configuration, dataset fingerprint, selected
  device, seed, and file paths.
\end{itemize}

Switch comparisons are stored separately from native model runs because the
same predictions can be converted into switch sequences under different
decision rules. Model-vs-reference comparisons are written under
\path{results/comparisons/switch_eval/<task_name>/<selection_id>/}. Cross-task
comparisons are written under
\path{results/comparisons/cross_task_switch/<selection_id>/}. This separation
keeps model training artifacts distinct from downstream operational
evaluation.

\section{Reproducibility Checklist}
\label{app:reproducibility-checklist}

The experiments are intended to be reproducible from configurations and saved
artifacts. The following information is recorded or fixed for each run:

\begin{itemize}
  \item the YAML configuration path and configuration fingerprint;
  \item the processed dataset path and dataset metadata fingerprint;
  \item the task formulation, input representation, target definition, context
  length, prediction horizon or duration horizon, and threshold assumptions;
  \item train, validation, and external test split metadata;
  \item random seed and hardware device used by the model runner;
  \item model hyperparameters, selected validation score, and final test
  metrics;
  \item saved predictions used for switch conversion;
  \item switch-conversion rule, post-processing settings, and Perfect Switch
  reference path.
\end{itemize}

The Python environment is specified by \path{environment.yml}. Experiments use
the active environment rather than creating a new virtual environment inside
the repository. GPU acceleration is used when available, with CUDA preferred,
then Apple MPS when supported, and CPU as a fallback. Grid searches distribute
independent trials over available CUDA devices when possible.
